
% HBL_INTERFACE_REGULARITY_GATE_20261111
\paragraph{Actual interface regularity must precede the energy-dual source estimate.}
Assume the one-sided original Hartree fields are classical, have
continuous time traces, and meet with $[U]=[U_r]=0$ along $r=L(s)>0$.
The moving trace identity gives $[U_s]=-L_s[U_r]=0$.
Where the actual nonlocal potential is continuous, the complete
renormalized physical residual satisfies
\[
 [F_{\rm full}]=[U_{rr}].
\]
Thus a nonzero curvature discrepancy in the literal four-piece
construction produces a jump in the actual source, not a Dirac
mass in the undifferentiated PDE. This distinction is important.
For $J=[F]_{r=R}\ne0$, choose a real compactly supported smooth
$\phi$ with $\phi'(0)=1$ and set
$g_\epsilon(r)=\sqrt\epsilon\,\phi((r-R)/\epsilon)$.
For the genuine phase block $L_-=-\Delta_r-3W^4$,
\[
 \|g_\epsilon\|_{\rm Mor}=O(1),\qquad
 (L_-g_\epsilon,F)_{L^2(\mathbb R^3)}
   =4\pi R^2J\epsilon^{-1/2}+O(\epsilon^{1/2}).
\]
Consequently the proposed Hessian--Morawetz-dual source norm is
infinite whenever this jump is nonzero. The bounded nonlocal
lower-order operator in $L_+$ does not remove the principal
obstruction. A nonzero jump also excludes local $H^{1/2}$
regularity, by the divergent integral of $(x+y)^{-2}$ across the
interface. A small positive time factor cannot make this infinite
per-time norm finite. These are conditional statements about
nonzero jumps, not a proof that every selected orbit has such a jump.

Merely repairing through $C^2$ is also insufficient for a general
moving-interface $H^1$-source Lipschitz assertion. Near $r=a$,
$U_a=(r-a)_+^3/6$ is $C^2$, but
$\partial_r\Delta_rU_a$ has a unit jump. After a smooth fixed-annulus
localization,
\[
 \|\Delta_r U_{a+h}-\Delta_r U_a\|_{H^1}\ge c\sqrt{|h|}.
\]
A finite repair through the third spatial jet removes this
particular source-derivative jump. With the \emph{actual} mismatch
jets $d_2=U_{\rm in}''-U_{\rm out}''$ and
$d_3=U_{\rm in}'''-U_{\rm out}'''$ at $L$, put
\[
 P(r)=h^2d_2\phi_2((r-L)/h)+h^3d_3\phi_3((r-L)/h),
\]
\[
 \phi_2(t)=\tfrac12t^2(1-t)^4(1+4t),\qquad
 \phi_3(t)=\tfrac16t^3(1-t)^4.
\]
Its endpoint jets are exactly
\[
 (P,P',P'',P''')(L)=(0,0,d_2,d_3),\qquad
 (P,P',P'',P''')(L+h)=(0,0,0,0).
\]
Only these polynomial identities and their factorizations are
assigned to the companion Lean module. The analytic obstruction
and the global physical splice are separate paper arguments.
For $k\le4$ the explicit source costs are
\[
 \|P^{(k)}\|_2\le
 h^{5/2-k}|d_2|\|\phi_2^{(k)}\|_2+
 h^{7/2-k}|d_3|\|\phi_3^{(k)}\|_2.
\]
No time-positive exponent follows from these formulas.
The full nonlocal Hartree source and the original two-speed
compatibility must be recomputed for the repaired field; preserving
boundary values and slopes alone does not preserve the convolution
or prove that the old speeds are still compatible. Actual moving
jet, fourth-derivative and parameter estimates remain open.
The complete proof and scope are recorded in
\texttt{HBL\_213\_INTERFACE\_REGULARITY\_GATE\_20261111.md}.
Original revised acceptance stays $11/13$.
